% Part: applied-modal-logics
% Chapter: temporal-logic
% Section: properties-accessibility

\documentclass[../../../include/open-logic-section]{subfiles}

\begin{document}

\olfileid{aml}{tl}{acc}

\olsection{સમયસંબંધી ફ્રેમોના ગુણધર્મો}

આપણાં સમયસંબંધી નિદર્શનોમાં સંબંધ~$\prec$ પર કોઈ શરત નથી,
એટલે આપણા પરિચિત સ્વયંસિદ્ધાંતોમાંથી બધાં નિદર્શનોમાં માન્ય
રહેતો એકમાત્ર સ્વયંસિદ્ધાંત~$K$, અથવા તેના અનુરૂપ
$K_{\Gtemp}$ અને~$K_{\Htemp}$ છે:
\begin{align*}
\tag{$K_{\Gtemp}$}  \Gtemp (p \to q) & \to (\Gtemp p \to \Gtemp q)\\
\tag{$K_{\Htemp}$}  \Htemp (p \to q) & \to (\Htemp p \to \Htemp q)
\end{align*}

પરંતુ સમયને કેવી રીતે દર્શાવવો તે અંગે નિદર્શનો પર વધુ કડક
શરતો મૂકીએ, દાખલા તરીકે~$\prec$ રેખીય ક્રમ હોય તે સુનિશ્ચિત
કરીએ, તો નિદર્શનોમાં બીજા દાવાઓ પણ માન્ય થશે.

\begin{table}[t]
  \begin{tabular}{| p{.48\textwidth} || p{.48\textwidth} |}
    \hline
    {\emph{જો $\prec$ \dots હોય}} & {\emph{તો~$\mModel{M}$માં \dots સત્ય છે:}} \\
    \hline \hline
    \emph{પરંપરિત}: \newline
    $\forall u \forall v \forall w ((u \prec v \land v \prec w) \lif u \prec w)$ & 
    $\Ftemp \Ftemp p \lif \Ftemp p$  \\
    \hline 
    \emph{રેખીય}: \newline
    $\forall w \forall v (w \prec v \lor w = v \lor v \prec w)$ &  
    $(\Ftemp \Ptemp  p \lor \Ptemp \Ftemp p) \lif (\Ptemp p \lor p \lor \Ftemp p)$\\
    \hline
    \emph{સઘન}: \newline
    $\forall w \forall v (w \prec v \to \exists u(w \prec u \land u \prec v))$ &  
    $\Ftemp p \lif \Ftemp \Ftemp p$ \\
    \hline
    \emph{ભૂતકાળ તરફ અસીમ}: \newline
    $\forall w \exists v( v \prec w)$ &  
    $\Htemp p \to \Ptemp p$ \\
    \hline
    \emph{ભવિષ્ય તરફ અસીમ}: \newline
    $\forall w \exists v( w \prec v)$ &  
    $\Gtemp p \to \Ftemp p$ \\
    \hline
  \end{tabular}
  \caption{સમયસંબંધી ફ્રેમોના કેટલાક અનુરૂપતા-ગુણધર્મો.}
  \ollabel{tab:correspondence}
\end{table} 

\olref{tab:correspondence}ના કેટલાક ગુણધર્મો સમય દર્શાવવા માટેના
નિદર્શનમાં ઇચ્છનીય લાગે. જોકે સંબંધ~$\prec$ પર આપણે ગમે તેવી
શરતો મૂકી શકીએ, તેમ છતાં બધી શરતોને સમયસંબંધી તર્કની ભાષામાં
વ્યક્ત થઈ શકતાં !!{formula}s અનુરૂપ હોતાં નથી. દાખલા તરીકે,
અસ્વવાચકતા, એટલે કે $\forall w \lnot (w \prec w)$ હોવાનો
ગુણધર્મ, સમયસંબંધી તર્કમાં કોઈ અનુરૂપ સૂત્ર ધરાવતો નથી.

\end{document}
